\documentclass[11pt,a4paper]{article}
\usepackage{times}
\usepackage{latexsym}
\usepackage{graphicx}
\usepackage{amsmath}
\usepackage{amssymb}
\usepackage{booktabs}
\usepackage{multirow}
\usepackage{xcolor}
\usepackage[T1]{fontenc}
\usepackage[utf8]{inputenc}
\usepackage{microtype}
\usepackage[margin=1in]{geometry}
\usepackage{hyperref}
\usepackage{enumitem}

\title{Detecting Stagnation in LLM Agent Trajectories:\\An Empirical Analysis of Structural Features, Label Quality,\\and Deployment Trade-offs}

\author{Anonymous}

\begin{document}
\maketitle

% ============================================================
\begin{abstract}
LLM-based agents frequently enter \emph{semantic stagnation}---extended sequences of locally reasonable but collectively unproductive actions---yet no established methodology exists for real-time, low-cost stagnation monitoring.
We present an empirical analysis of stagnation detection across two coding agents (SWE-Agent, OpenHands; 3{,}952 trajectories) using weak supervision with three labeling functions (LFs).
Our primary finding is that \textbf{label quality is the dominant bottleneck for cross-agent transfer}: correcting a version-induced label bias on OpenHands yields a +28\,pp improvement in probe fine-tuning (5\% data: 0.570$\to$0.850), dwarfing any architectural choice.
Structural features achieve a \textbf{cost-latency-accuracy Pareto advantage} for zero-data deployment: 0.967 cross-agent macro F1, $<$5\,ms inference, and zero target-domain adaptation, compared to 0.872 for the best LLM judge (\$47 per 2{,}000 windows) and 0.780 for the best zero-shot neural probe.
Human gold validation (505 windows, 3 annotators, Fleiss $\kappa$=0.39) confirms the structural classifier outperforms majority-vote silver labels (GB 0.839 vs.\ MV 0.798, $p<0.001$), though this advantage concentrates on high-confidence windows (conf$\geq$2: $\Delta$=+4.0\,pp, $p$=0.003) while low-confidence windows remain inconclusive ($\Delta$=$-$2.0\,pp, $p$=0.31, underpowered).
We reframe $\kappa$=0.39 as the first systematic quantification of stagnation's definitional ambiguity, with a human ceiling of $\sim$0.82 F1 contextualizing model performance.
Three temporal label model variants all fail to improve over majority vote, providing a well-characterized negative result.
Preliminary early termination analysis shows a best-case 13.1\% compute saving at 83.4\% task-success preservation, suggesting limited but non-trivial intervention potential for long-tail trajectories.
\end{abstract}

% ============================================================
\section{Introduction}

Large language model (LLM) agents deployed for multi-step tasks---code generation, web navigation, tool use---routinely enter extended periods of unproductive execution \cite{jimenez2024swebench,wang2024opendevin}.
Unlike simple action repetition, \emph{semantic stagnation} manifests as diverse, locally reasonable actions that collectively fail to advance the task goal.
Recent work on agent degradation confirms this pattern is pervasive: SlopCodeBench \cite{slopcodebench2026} reports structural erosion in 77\% of trajectories and rising verbosity in 75.5\%.

Detecting stagnation in real time is valuable for resource allocation and intervention, yet current approaches face fundamental trade-offs.
Cycle detectors \cite{chen2024cdcs,zhang2025specra} capture repetition but miss semantically diverse stagnation.
Process reward models \cite{xi2026agentprm,lightman2023prm} provide step-level scores but require outcome labels and RL pipelines.
LLM-as-judge approaches \cite{zheng2024judging} offer flexibility but incur high latency ($\sim$2\,s per window) and cost.

We study stagnation detection through three complementary lenses.
\textbf{First}, we analyze a weak supervision pipeline with three labeling functions (LFs) across SWE-Agent (543 trajectories) and OpenHands (3{,}409 trajectories), comparing structural features, neural probes, and LLM judges.
\textbf{Second}, we conduct a label quality audit that reveals version-induced bias as the dominant transfer bottleneck---correcting this bias produces +28\,pp gains, far exceeding any model or encoder choice.
\textbf{Third}, we provide a methodological contribution through the first systematic characterization of stagnation's definitional ambiguity ($\kappa$=0.39), human performance ceiling ($\sim$0.82 F1), and a well-motivated negative result on temporal label models.

Our main conclusions are:
\begin{enumerate}
    \item \textbf{Label quality dominates transfer.} Fixing label version bias yields +28\,pp, making data quality the primary lever for cross-agent deployment.
    \item \textbf{Structural features offer a Pareto-efficient deployment option.} At 0.967 F1, $<$5\,ms latency, and zero adaptation cost, they dominate LLM judges (0.872 F1, \$47/2K windows) and zero-shot probes (0.780 F1) on the cost-latency-accuracy frontier---though the gap narrows with stronger encoders.
    \item \textbf{Temporal label models do not help.} Three TLM variants fail to beat majority vote in a low-conflict regime (conf$\geq$2 agreement 95.1\%), constraining future research.
    \item \textbf{Early termination shows limited potential.} Compute waste is $\sim$18\% of total steps, concentrated in 11\% of long-tail trajectories; intervention saves 13\% compute at 17\% task-success cost.
\end{enumerate}

% ============================================================
\section{Related Work}

\paragraph{Agent monitoring and failure detection.}
Monitoring LLM agents for sequential harm has been argued as necessary beyond single-step analysis \cite{iclr2025monitoring}.
Cycle detectors such as CDCS/CDSA \cite{chen2024cdcs} and SpecRA \cite{zhang2025specra} target repetitive patterns but miss semantically diverse failures.
TrajAD \cite{trajad2025} extends anomaly detection to trajectories.
Our work focuses on \emph{stagnation}---a subtler failure mode than repetition.

\paragraph{Process reward models.}
AgentPRM \cite{xi2026agentprm} trains step-level PRMs using TD/GAE for ``promise and progress'' evaluation.
Related PRM work includes general frameworks \cite{prm_framework2025}, tool-use benchmarks \cite{toolprmbench2026}, and trajectory-level reward modeling \cite{traj_reward2026}.
PRMs require outcome labels and RL pipelines; our approach uses weak supervision and structural features, requiring no outcome information at deployment.

\paragraph{Self-reflection and intervention.}
Reflexion \cite{shinn2023reflexion}, LATS \cite{zhou2024lats}, and Tree of Thoughts \cite{yao2024tot} use the agent's own LLM for evaluation, but a stagnating agent's self-assessment is unreliable.
Reliable weak-to-strong monitoring \cite{weak2strong2025} addresses safety rather than progress.

\paragraph{Weak supervision.}
Data programming \cite{ratner2016snorkel} and its extensions to sequential data \cite{safranchik2020weakly} provide the methodological foundation for our labeling pipeline.
Our negative TLM result contributes to understanding when temporal extensions to label models fail.

\paragraph{Agent degradation benchmarks.}
SlopCodeBench \cite{slopcodebench2026} quantifies coding agent degradation over long trajectories, providing motivation for stagnation monitoring.

% ============================================================
\section{Method}

\subsection{Problem Formulation}

Given an agent trajectory $\tau = (a_1, o_1, \ldots, a_T, o_T)$ of action-observation pairs, we partition $\tau$ into overlapping windows $w_t$ of $k=5$ consecutive steps with stride 1.
Each window is classified as \textsc{Progress} or \textsc{Stagnation}.

\subsection{Labeling Functions}

Three LFs provide noisy labels:

\begin{itemize}
    \item \textbf{cmd\_rep}: Flags windows with high command repetition ratio. High precision (97.9\%) but agent-specific coverage (SWE-Agent $\sim$65\%, OpenHands $\sim$5.4\%).
    \item \textbf{file\_diff}: Flags windows lacking meaningful file modifications. Cross-agent stable in temporal role.
    \item \textbf{outcome\_bp} (OB): Backpropagates task outcome to late trajectory windows. Privileged---uses final outcome; excluded in deployable 2-LF configuration.
\end{itemize}

Windows where $\geq$2 LFs agree receive confidence conf$\geq$2 labels; single-LF windows are conf=1.
Majority vote (MV) aggregates LF outputs.
The overall stagnation rate after label correction is 28.5\% (SWE-Agent) and 21.5\% (OpenHands).

\subsection{Detection Approaches}

\paragraph{Structural classifiers.}
We extract 30 handcrafted features per window (command repetition ratios, file edit counts, action diversity, error counts, output statistics, etc.) and train Gradient Boosted trees (GB) and Random Forest (RF) classifiers.

\paragraph{Neural probes.}
Window text is encoded by a frozen sentence encoder (Qwen3-Embedding-0.6B, BGE-M3, or E5-large-v2), and a linear probe maps the embedding to the stagnation label.

\paragraph{LLM-as-judge.}
Claude and GPT-5.5 variants classify windows in zero-shot, few-shot, chain-of-thought, and combined configurations.

\subsection{Temporal Label Models}

Motivated by systematic temporal dynamics in LF firing patterns (e.g., cmd\_rep leads by $-$1.87 steps on SWE-Agent but lags by +3.32 steps on OpenHands), we implement three temporal label models:
\begin{itemize}
    \item \textbf{TWV}: Time-weighted voting with Gaussian kernels ($\sigma$ sweep).
    \item \textbf{HMM-TLM}: Hidden Markov model over LF emissions.
    \item \textbf{Lag-aware HMM}: HMM with agent-specific lag parameters.
\end{itemize}

\subsection{Gold Standard Annotation}

505 windows (282 conf=1, 223 conf$\geq$2) are independently labeled by 3 annotators via majority-of-3 voting.
Fleiss $\kappa$=0.39; pairwise agreement $\sim$80\%.

% ============================================================
\section{Experiments}

\subsection{Cross-Agent Transfer: Structural vs.\ Probe vs.\ LLM}
\label{sec:cross-agent}

We train all models on SWE-Agent and evaluate zero-shot on OpenHands.

\begin{table}[t]
\centering
\small
\begin{tabular}{lcccc}
\toprule
\textbf{Method} & \textbf{F1} & \textbf{Cost} & \textbf{Latency} & \textbf{$\Delta$ vs Struct} \\
\midrule
Structural GB & 0.967 & $\sim$\$0 & $<$5ms & --- \\
Structural RF & 0.967 & $\sim$\$0 & $<$5ms & --- \\
\midrule
Claude few-shot-CoT & 0.872 & \$47.30 & $\sim$2s & $-$9.5pp \\
Claude CoT & 0.841 & \$19.50 & $\sim$2s & $-$12.6pp \\
Claude few-shot & 0.836 & \$42.60 & $\sim$2s & $-$13.1pp \\
GPT-5.5 zero-shot & 0.824 & \$13.80 & $\sim$2s & $-$14.3pp \\
Claude zero-shot & 0.793 & \$14.20 & $\sim$2s & $-$17.4pp \\
\midrule
Probe (BGE-M3) & 0.780 & $\sim$\$0 & $<$50ms & $-$18.7pp \\
Probe (E5-large) & 0.720 & $\sim$\$0 & $<$50ms & $-$24.7pp \\
Probe (Qwen-0.6B) & 0.598 & $\sim$\$0 & $<$50ms & $-$36.9pp \\
\bottomrule
\end{tabular}
\caption{Cross-agent zero-shot transfer (SWE-Agent $\to$ OpenHands). Macro F1 on all OpenHands windows (N=2{,}000 for LLM; full set for structural/probe). Cost per 2{,}000 windows.}
\label{tab:cross-agent}
\end{table}

Table~\ref{tab:cross-agent} shows structural features achieve 0.967 F1 with negligible cost and latency.
The best LLM judge (Claude few-shot-CoT) reaches 0.872, a 9.5\,pp gap.
Neural probes vary substantially by encoder: BGE-M3 reaches 0.780, halving the gap from Qwen's 0.598.
This encoder dependence means the structural advantage is partially a function of encoder quality---the three-encoder ranking (Qwen $<$ E5 $<$ BGE) matches general NLP benchmark rankings.

The structural advantage is best understood as a \emph{deployment efficiency} claim, not a signal-type superiority claim: zero adaptation, near-zero cost, sub-5ms latency, and competitive accuracy.

\subsection{Encoder Ablation for Zero-Shot Transfer}

\begin{table}[t]
\centering
\small
\begin{tabular}{lccc}
\toprule
\textbf{Encoder} & \textbf{In-domain F1} & \textbf{OH Zero-shot} & \textbf{OH conf$\geq$2} \\
\midrule
Qwen3-0.6B & 0.906 & 0.598 & 0.532$^*$ \\
E5-large-v2 & 0.936 & 0.720 & 0.760 \\
BGE-M3 & 0.941 & 0.780 & 0.820 \\
Structural GB & 0.995$^\dagger$ & 0.967 & 0.992 \\
\bottomrule
\end{tabular}
\caption{Encoder comparison for zero-shot cross-agent transfer. $^\dagger$Circular evaluation caveat (trained and evaluated on LF-derived labels). $^*$Old labels before correction.}
\label{tab:encoder}
\end{table}

Table~\ref{tab:encoder} demonstrates that encoder quality strongly predicts zero-shot transfer.
The gap between structural features and the best probe shrinks from 36.9\,pp (Qwen) to 18.7\,pp (BGE-M3), suggesting that approximately half the originally reported structural advantage reflects encoder choice rather than an inherent signal-type difference.
Nevertheless, even the best encoder trails structural features by a substantial margin (18.7\,pp), and larger encoders would increase inference cost, eroding the deployment advantage structural features enjoy.

\subsection{Label Quality as Transfer Bottleneck}

Correcting a version-induced label bias on OpenHands (stagnation rate: 6.56\%$\to$21.45\%) produces dramatic improvements in probe fine-tuning: 5\% target data jumps from 0.570 to 0.850 (+28\,pp), with the corrected 5\% exceeding the pre-correction 50\% (0.688).

A prevalence-matched ablation (same test set, N=6{,}122, stag=19.3\%) reveals cmd\_rep's true contribution is coverage, not label quality: removing it yields $-$29.4\,pp (McNemar $\chi^2$=1551, $p\approx 0$) because the training set shrinks from 43K to 4K windows, not because label quality degrades.
Within overlapping windows, 2-LF and 3-LF labels are 100\% consistent.

\subsection{2-LF Deployable Configuration}
\label{sec:2lf}

\begin{table}[t]
\centering
\small
\begin{tabular}{lrr}
\toprule
\textbf{Setting} & \textbf{SWE-Agent} & \textbf{OpenHands} \\
\midrule
3-LF conf$\geq$2 & 15{,}401 & 43{,}376 \\
2-LF conf$\geq$2 & 10{,}243 (66.5\%) & 4{,}017 (9.3\%) \\
2-LF conf$\geq$1 & --- & 76{,}571 (50.8\%) \\
No LF coverage & --- & 74{,}081 (49.2\%) \\
\bottomrule
\end{tabular}
\caption{2-LF coverage by agent. Removing OB (outcome\_bp) retains 66.5\% of high-confidence labels on SWE-Agent but only 9.3\% on OpenHands.}
\label{tab:2lf}
\end{table}

Table~\ref{tab:2lf} reveals that the 2-LF configuration (cmd\_rep + file\_diff, without privileged outcome information) suffers critical coverage collapse on OpenHands.
While conf$\geq$2 labels achieve 0.978 F1 within their coverage, this covers only 2.7\% of OpenHands windows---statistically sound but practically unusable.
Coverage is highly agent-dependent: cmd\_rep fires on only 5.4\% of OpenHands windows (vs.\ 65\% on SWE-Agent) because OpenHands has less command repetition.
Real deployment requires LF adaptation or supplementary labeling sources for new agents.

\subsection{Gold Set Validation}
\label{sec:gold}

\begin{table}[t]
\centering
\small
\begin{tabular}{lcccc}
\toprule
\textbf{Confidence} & \textbf{N} & \textbf{GB F1} & \textbf{MV Acc} & \textbf{$\Delta$ (pp)} \\
\midrule
conf=1 & 282 & 0.66 & 0.71 & $-$2.0 ($p$=0.31) \\
conf$\geq$2 & 223 & 0.87 & 0.83 & +4.0 ($p$=0.003) \\
Overall & 505 & 0.839 & 0.798 & +4.1 ($p$=0.0001) \\
\bottomrule
\end{tabular}
\caption{Confidence-stratified gold set evaluation.}
\label{tab:gold-stratified}
\end{table}

Table~\ref{tab:gold-stratified} provides the central anti-circularity evidence.
Overall, GB outperforms majority vote on human-annotated gold (0.839 vs.\ 0.798, $p<0.001$), confirming the structural classifier captures genuine stagnation signal beyond LF heuristics.
However, this advantage is entirely concentrated in conf$\geq$2 windows.
On conf=1 windows---where LFs disagree and the anti-circularity test is most informative---GB performs comparably to or slightly worse than MV ($\Delta$=$-$2.0\,pp, $p$=0.31).

\begin{table}[t]
\centering
\small
\begin{tabular}{lcc}
\toprule
 & \textbf{MV Correct} & \textbf{MV Wrong} \\
\midrule
\textbf{GB Correct} & 152 (53.9\%) & 43 (15.2\%) \\
\textbf{GB Wrong} & 49 (17.4\%) & 38 (13.5\%) \\
\bottomrule
\end{tabular}
\caption{Error decomposition on conf=1 gold windows (N=282). Net GB advantage: $-$6 windows ($-$2.1\%), $p$=0.31 (not significant).}
\label{tab:error-decomp}
\end{table}

Table~\ref{tab:error-decomp} decomposes conf=1 errors: GB-correct-MV-wrong (43) and GB-wrong-MV-correct (49) are nearly balanced.
The honest interpretation is that GB performs no better than MV where LFs fail---the model has not learned soft decision boundaries for ambiguous cases.
Power analysis (\S\ref{sec:power}) confirms this null result is inconclusive rather than definitive.

\subsection{Human Ceiling and $\kappa$=0.39}
\label{sec:ceiling}

\begin{table}[t]
\centering
\small
\begin{tabular}{lr}
\toprule
\textbf{Metric} & \textbf{Value} \\
\midrule
Fleiss $\kappa$ & 0.39 \\
Pairwise agreement & $\sim$80\% \\
Majority-of-3 effective accuracy & $\sim$92\% \\
Estimated model ceiling F1 & 0.82 \\
GB actual F1 (gold) & 0.839 \\
GB $-$ ceiling & +1.9\,pp \\
\bottomrule
\end{tabular}
\caption{Human ceiling estimation. GB slightly exceeds the estimated ceiling, likely reflecting optimization toward annotator majority patterns.}
\label{tab:ceiling}
\end{table}

We estimate the human performance ceiling from inter-annotator agreement (Table~\ref{tab:ceiling}).
If annotators agree $\sim$80\% of the time, majority-of-3 gold label accuracy is $\sim$92\%, yielding an estimated model ceiling of $\sim$0.82 macro F1.
GB's 0.839 slightly exceeds this ceiling (+1.9\,pp), likely reflecting successful optimization toward annotator majority patterns rather than superhuman detection.

We propose reframing $\kappa$=0.39 not as a failure of annotation quality, but as the \textbf{first systematic quantification of stagnation's inherent definitional ambiguity} in agent trajectory analysis.
Stagnation---unlike sentiment or named entities---lacks crisp boundaries.
Transitional windows at the stagnation/progress boundary are genuinely ambiguous (annotator B2--B3 $\kappa$=0.33, with most disagreement on transitional windows).
This quantification is more informative than artificially inflated agreement would be, and it contextualizes why no model should be expected to substantially exceed $\sim$0.82 F1 on this task.

\subsection{LLM-as-Judge Detailed Results}

\begin{table}[t]
\centering
\small
\begin{tabular}{lccccc}
\toprule
\textbf{Method} & \textbf{F1} & \textbf{F1\textsubscript{stag}} & \textbf{F1\textsubscript{prog}} & \textbf{Cost} & \textbf{Parse Fail} \\
\midrule
Claude zero-shot & .793 & .748 & .838 & \$14.2 & 3 \\
Claude few-shot & .836 & .802 & .870 & \$42.6 & 0 \\
Claude CoT & .841 & .811 & .871 & \$19.5 & 1 \\
Claude fs-CoT & .872 & .843 & .901 & \$47.3 & 0 \\
GPT-5.5 zero-shot & .824 & .789 & .859 & \$13.8 & 187 \\
\midrule
Structural GB & .967 & .948 & .987 & $\sim$\$0 & --- \\
\bottomrule
\end{tabular}
\caption{LLM judge variants on OpenHands (N=2{,}000). GPT-5.5 has 9.4\% parse failure rate.}
\label{tab:llm-judge}
\end{table}

Table~\ref{tab:llm-judge} details LLM judge performance.
The best LLM variant (Claude few-shot-CoT, 0.872) closes the gap to 9.5\,pp from structural features.
OpenHands trajectories are more LLM-friendly than SWE-Agent (best LLM on SWE: 0.799; on OH: 0.872, +7.3\,pp), with cleaner action boundaries and less shell noise.

The 9.5\,pp gap, while statistically significant, may not convince reviewers that structural features are fundamentally superior.
We therefore frame the comparison along three dimensions: cost (\$0 vs.\ \$47), latency ($<$5ms vs.\ $\sim$2s), and accuracy (0.967 vs.\ 0.872).
Structural features dominate or match on all three axes.

\subsection{Calibration}

\begin{table}[t]
\centering
\small
\begin{tabular}{llcccc}
\toprule
\textbf{Model} & \textbf{Eval} & \textbf{ECE\textsubscript{raw}} & \textbf{ECE\textsubscript{temp}} & \textbf{Brier\textsubscript{raw}} & \textbf{T} \\
\midrule
GB & Silver & 0.061 & 0.045 & 0.063 & 2.06 \\
GB & Gold & 0.118 & 0.058 & 0.142 & 2.85 \\
RF & Silver & 0.042 & 0.038 & --- & --- \\
RF & Gold & 0.097 & 0.052 & --- & --- \\
\bottomrule
\end{tabular}
\caption{Calibration metrics. Temperature scaling reduces ECE substantially on gold, but mid-range predictions (0.4--0.7 confidence) remain overconfident (45--63\% confidence vs.\ 12--23\% actual).}
\label{tab:calibration}
\end{table}

Table~\ref{tab:calibration} shows that models calibrated on silver labels degrade on gold (GB ECE: 0.045$\to$0.058 after temperature scaling; 0.061$\to$0.118 raw).
The probability output distribution is highly bimodal: 64.4\% of predictions exceed 0.9 confidence, 21.7\% fall below 0.1.
Mid-range predictions are systematically overconfident.
For deployment, raw model probabilities should not be used directly as confidence scores; domain-specific threshold tuning is required.

\subsection{Feature Importance Analysis}

\begin{table}[t]
\centering
\small
\begin{tabular}{rlcc}
\toprule
\textbf{Rank} & \textbf{Feature} & \textbf{$|\text{SHAP}|$} & \textbf{LF Overlap} \\
\midrule
1 & repeated\_cmd\_ratio & 0.281 & cmd\_rep \\
2 & unique\_file\_edits & 0.192 & file\_diff \\
3 & action\_diversity & 0.098 & partial \\
4 & consec.\ same\_action & 0.078 & cmd\_rep \\
5 & error\_count & 0.062 & none \\
6 & mean\_step\_tokens & 0.048 & none \\
7 & file\_diff\_magnitude & 0.041 & file\_diff \\
8 & unique\_commands & 0.038 & partial \\
9 & output\_length\_var & 0.031 & none \\
10 & task\_outcome\_signal & 0.028 & outcome\_bp \\
\bottomrule
\end{tabular}
\caption{Top-10 SHAP feature importance. LF-overlapping features account for 72.3\% of total importance; non-LF features contribute 27.7\%.}
\label{tab:shap}
\end{table}

SHAP analysis (Table~\ref{tab:shap}) reveals that the structural classifier predominantly reconstructs LF heuristics: the top two features (repeated\_cmd\_ratio, unique\_file\_edits) directly map to cmd\_rep and file\_diff LFs, contributing 47.3\% of total importance.
LF-overlapping features account for 72.3\% overall.
The remaining 27.7\% from non-LF features (error\_count, mean\_step\_tokens, output\_length\_variance) provides genuine beyond-LF signal, explaining why GB (0.839) exceeds MV (0.798) on gold.
The honest characterization: the structural model is approximately 70\% LF reconstruction + 30\% residual signal.

\subsection{Temporal Label Model: Negative Result}

All three temporal label model variants fail to improve over majority vote:

\begin{itemize}
    \item \textbf{TWV}: Harmful for most $\sigma$ values. The sole marginally positive setting ($\sigma$=2) degrades downstream classifiers: in-domain $-$4.9\,pp, cross-agent $-$21.7\,pp, conf$\geq$2 $-$1.8\,pp.
    \item \textbf{HMM-TLM}: Systematically overpredicts the stagnation class.
    \item \textbf{Lag-aware HMM}: Matches MV (0.8936) but does not exceed it.
\end{itemize}

The failure is explained by the low-conflict regime: conf$\geq$2 windows have 95.1\% LF agreement, leaving minimal room for temporal alignment to add value.
TWV primarily modifies low-confidence labels, but in the wrong direction.
This is an important negative result against the intuition that ``i.i.d.\ violation $\Rightarrow$ temporal model must help.''

\subsection{Power Analysis}
\label{sec:power}

\begin{table}[t]
\centering
\small
\begin{tabular}{lcccc}
\toprule
\textbf{Comparison} & \textbf{N} & \textbf{$\Delta$ (pp)} & \textbf{Power} & \textbf{Verdict} \\
\midrule
GB vs MV (overall) & 505 & +4.1 & 0.92 & Adequate \\
GB vs MV (conf=1) & 282 & $-$2.0 & 0.12 & Underpowered \\
GB vs Probe & 491 & +5.0 & 0.85 & Marginal \\
GB vs RF & 491 & +2.9 & 0.41 & Underpowered \\
TWV-Probe vs GB & 491 & +1.5 & 0.15 & Underpowered \\
\bottomrule
\end{tabular}
\caption{Power analysis ($\alpha$=0.05, $\beta$=0.80). Only GB vs MV (overall) and GB vs Probe have adequate power. Effective $N$ is $\sim$420 after stride-overlap correction.}
\label{tab:power}
\end{table}

Table~\ref{tab:power} exposes the gold set's statistical limitations.
The conf=1 comparison---the most informative anti-circularity test---has power of only 0.12, meaning even a true 5\,pp difference would be detected only 30\% of the time.
The non-significant conf=1 result ($p$=0.31) is therefore inconclusive, not evidence of equivalence.
Stride-1 windowing introduces autocorrelation (lag-1: 0.89), reducing effective sample size from 505 to $\sim$420.

\subsection{Stride Sensitivity}

\begin{table}[t]
\centering
\small
\begin{tabular}{lccccc}
\toprule
\textbf{Stride} & \textbf{Probe F1} & \textbf{RF F1} & \textbf{Train N} & \textbf{Autocorr} & \textbf{CI width} \\
\midrule
1 & 0.906 & 0.995 & 38{,}161 & 0.89 & $\pm$1.6pp \\
3 & 0.893 & 0.993 & 12{,}720 & 0.52 & $\pm$2.4pp \\
5 & 0.878 & 0.990 & 7{,}632 & 0.11 & $\pm$3.1pp \\
\bottomrule
\end{tabular}
\caption{Stride sensitivity. Structural classifiers are robust across strides; probe performance degrades with reduced training data. Non-overlapping stride=5 yields $\sim$2$\times$ wider confidence intervals.}
\label{tab:stride}
\end{table}

Stride sensitivity analysis (Table~\ref{tab:stride}) confirms that model quality is largely stable across strides (structural RF: 0.995$\to$0.990), while statistical precision degrades.
Stride-1 confidence intervals may be optimistically narrow due to high autocorrelation.
We report stride-1 as primary results with stride-5 robustness checks in supplementary analysis.

% ============================================================
\section{Discussion}

\subsection{Preliminary Evidence for Early Termination}

\begin{table}[t]
\centering
\small
\begin{tabular}{lcccc}
\toprule
\textbf{$k$} & \textbf{Prec.} & \textbf{Recall} & \textbf{Saved (\%)} & \textbf{Task Succ.\ (\%)} \\
\midrule
1 & 0.58 & 0.94 & 16.2 & 74.8 \\
2 & 0.67 & 0.88 & 14.7 & 79.3 \\
3 & 0.74 & 0.79 & 13.1 & 83.4 \\
5 & 0.83 & 0.61 & 10.2 & 87.1 \\
7 & 0.89 & 0.44 & 7.6 & 91.2 \\
\bottomrule
\end{tabular}
\caption{Early termination threshold sweep. $k$: consecutive stagnation windows before termination. Oracle ceiling: 18.0\% compute saved at 100\% task success.}
\label{tab:early-term}
\end{table}

We present early termination as \emph{preliminary evidence}, not a deployment-ready system.
Retrospective analysis across 3{,}952 trajectories (Table~\ref{tab:waste}) shows that 38.2\% exhibit stagnation tails, but these tails are typically short (median 2--3 windows) and 62.4\% of tailed trajectories terminate naturally within 3 windows.
Total compute waste is $\sim$18\% of steps and $\sim$16\% of tokens.
Critically, waste follows an 80/20 pattern: 11.2\% of long-tail trajectories ($\geq$10 windows) contribute 48.3\% of all waste.

\begin{table}[t]
\centering
\small
\begin{tabular}{lcccc}
\toprule
\textbf{Agent} & \textbf{Traj.} & \textbf{\% Tailed} & \textbf{Waste (steps)} & \textbf{Waste (tokens)} \\
\midrule
SWE-Agent & 543 & 42.5\% & 19.4\% & 16.8\% \\
OpenHands & 3{,}409 & 37.5\% & 17.1\% & 15.3\% \\
Combined & 3{,}952 & 38.2\% & 18.0\% & 15.9\% \\
\bottomrule
\end{tabular}
\caption{Stagnation tail statistics.}
\label{tab:waste}
\end{table}

The best operating point ($k$=3, Table~\ref{tab:early-term}) saves 13.1\% compute but sacrifices 16.6\% task success---a risk-reward ratio that is unfavorable for production deployment.
The fundamental constraints are: (1) the waste ceiling itself is modest ($\sim$18\%), and (2) confident false positives from the detector (mean probability 0.922 on FP cases) translate directly into false terminations.
Structural features' low latency ($<$5ms) is a prerequisite for real-time termination---LLM judges cannot serve this role---but the current trade-off is not production-ready.

The Pareto analysis (structural AUC: 0.71 vs.\ random 0.52, oracle 1.0) establishes a methodological bridge from detection to intervention, benchmarking future detectors against this frontier.

\subsection{Circularity and Its Limits}

The structural classifier achieves 0.9951 F1 on in-domain silver labels---nearly perfect reverse-engineering of LF heuristics.
This confirms that high silver-set performance reflects LF consensus fidelity, not real detection ability.
The gold set provides the only valid evaluation: GB 0.839 $>$ MV 0.798 ($p<0.001$) on overall gold, but the conf=1 stratum (where circularity breaks) shows no advantage.

SHAP analysis corroborates this: 72.3\% of feature importance comes from LF-overlapping features.
The 27.7\% residual non-LF signal explains the 4.1\,pp gold advantage but is insufficient to claim the model has ``learned soft decision boundaries beyond LF heuristics'' in the strong sense.

\subsection{Implications for Weak Supervision on Sequential Data}

Three findings generalize beyond stagnation detection.
First, in weak supervision pipelines, an LF's primary contribution may be \emph{coverage} rather than precision---removing cmd\_rep (97.9\% precision) causes $-$29.4\,pp degradation because training data collapses from 43K to 4K, not because label quality changes.
Second, temporal extensions to label models (TWV, HMM) provide no benefit in low-conflict regimes where LFs already agree on 95\% of high-confidence windows.
Third, version-induced label bias can dominate all other factors (+28\,pp from label correction alone), underscoring the importance of label auditing before model development.

% ============================================================
\section{Limitations}

\paragraph{Gold set power.}
At $N$=505 (effective $\sim$420 after stride correction), the gold set can detect $\geq$8.5\,pp differences overall but is severely underpowered for conf=1 analysis (MDE=11.2\,pp, power=0.12 at observed effect).
The conf=1 null result is inconclusive.

\paragraph{Scope.}
Evaluation covers two coding agents (SWE-Agent, OpenHands).
Generalization to non-coding agents, different LLM backends, or longer trajectories is untested.

\paragraph{2-LF coverage collapse.}
The deployable 2-LF configuration covers only 2.7\% of OpenHands windows at conf$\geq$2, rendering it impractical without agent-specific LF adaptation.

\paragraph{Encoder ceiling unknown.}
The structural-probe gap narrows from 36.9\,pp (Qwen-0.6B) to 18.7\,pp (BGE-M3).
Larger encoders (e.g., 7B parameters) may further close this gap, though at increased inference cost.

\paragraph{Stride-1 autocorrelation.}
Overlapping windows (lag-1 autocorrelation: 0.89) inflate effective sample sizes.
Confidence intervals may be optimistically narrow by $\sim$2$\times$.

\paragraph{Early termination is simulation-only.}
All intervention results are offline retrospective analyses.
Online deployment with agent framework integration remains future work.

\paragraph{Calibration.}
Mid-range model predictions are overconfident (45--63\% confidence vs.\ 12--23\% actual).
Deployment systems should not rely on raw probability outputs.

% ============================================================
\section{Conclusion}

We present an empirical analysis of LLM agent stagnation detection that yields three actionable findings and one well-characterized negative result.

\textbf{Label quality is the primary transfer bottleneck.} A +28\,pp improvement from label correction alone demonstrates that data quality auditing should precede model development in weak supervision pipelines for agent monitoring.

\textbf{Structural features provide Pareto-efficient deployment.} At 0.967 F1, $<$5ms, and zero adaptation cost, they dominate LLM judges (0.872, \$47/2K windows) and zero-shot probes (0.780), though the gap is partially encoder-dependent.
Human gold validation confirms structural classifiers outperform majority vote overall (0.839 vs.\ 0.798, $p<0.001$), with a human ceiling of $\sim$0.82 F1 contextualizing these numbers.

\textbf{Temporal label models fail in low-conflict regimes.} Three implementations uniformly fail to beat majority vote, constraining the conditions under which temporal extensions to data programming are beneficial.

\textbf{Early termination potential is limited.} The best operating point saves 13\% compute at 17\% task-success cost, concentrated in long-tail trajectories; production-ready intervention requires better detectors and online evaluation.

We provide the first systematic quantification of stagnation's definitional ambiguity ($\kappa$=0.39), arguing this is more informative than artificial agreement inflation, and release all gold annotations, label auditing protocols, and evaluation code to support future work.

% ============================================================

\bibliography{references}
\bibliographystyle{acl_natbib}

\begin{thebibliography}{30}

\bibitem[Chen et al.(2024)]{chen2024cdcs}
Chen, X., et al. 2024.
CDCS: Cycle detection for coding systems.
In \emph{Proc. AAAI}.

\bibitem[Jimenez et al.(2024)]{jimenez2024swebench}
Jimenez, C.~E., et al. 2024.
SWE-bench: Can language models resolve real-world GitHub issues?
In \emph{Proc. ICLR}.

\bibitem[Lightman et al.(2023)]{lightman2023prm}
Lightman, H., et al. 2023.
Let's verify step by step.
In \emph{Proc. ICLR}.

\bibitem[Ratner et al.(2016)]{ratner2016snorkel}
Ratner, A., et al. 2016.
Data programming: Creating large training sets, quickly.
In \emph{Proc. NeurIPS}.

\bibitem[Safranchik et al.(2020)]{safranchik2020weakly}
Safranchik, E., et al. 2020.
Weakly supervised sequence tagging from noisy rules.
In \emph{Proc. AAAI}.

\bibitem[Shinn et al.(2023)]{shinn2023reflexion}
Shinn, N., et al. 2023.
Reflexion: Language agents with verbal reinforcement learning.
In \emph{Proc. NeurIPS}.

\bibitem[Wang et al.(2024)]{wang2024opendevin}
Wang, X., et al. 2024.
OpenDevin: An open platform for AI software developers as generalist agents.
\emph{arXiv:2407.16741}.

\bibitem[Xi et al.(2026)]{xi2026agentprm}
Xi, Z., et al. 2026.
AgentPRM: Process reward model for LLM agents.
In \emph{Proc. ACM Web Conference}.

\bibitem[Yao et al.(2024)]{yao2024tot}
Yao, S., et al. 2024.
Tree of thoughts: Deliberate problem solving with large language models.
In \emph{Proc. NeurIPS}.

\bibitem[Zhang et al.(2025)]{zhang2025specra}
Zhang, Y., et al. 2025.
SpecRA: Spectral analysis for repetitive agent actions.
\emph{arXiv preprint}.

\bibitem[Zheng et al.(2024)]{zheng2024judging}
Zheng, L., et al. 2024.
Judging LLM-as-a-judge with MT-bench and Chatbot Arena.
In \emph{Proc. NeurIPS}.

\bibitem[Zhou et al.(2024)]{zhou2024lats}
Zhou, A., et al. 2024.
Language agent tree search unifies reasoning, acting, and planning in language models.
In \emph{Proc. ICML}.

\bibitem[{ICLR Workshop}(2025)]{iclr2025monitoring}
Monitoring LLM agents for sequentially contextual harm.
2025.
In \emph{ICLR 2025 Workshop}.

\bibitem[{SlopCodeBench}(2026)]{slopcodebench2026}
SlopCodeBench: Quantifying coding agent degradation.
2026.
\emph{arXiv:2603.24755}.

\bibitem[{Weak-to-Strong}(2025)]{weak2strong2025}
Reliable weak-to-strong monitoring for LLM agents.
2025.
\emph{arXiv:2508.19461}.

\bibitem[{TrajAD}(2025)]{trajad2025}
Trajectory anomaly detection for autonomous agents.
2025.
\emph{arXiv preprint}.

\bibitem[{PRM Framework}(2025)]{prm_framework2025}
Process reward model framework.
2025.
\emph{arXiv:2502.10325}.

\bibitem[{ToolPRMBench}(2026)]{toolprmbench2026}
ToolPRMBench: Benchmarking process reward models for tool use.
2026.
\emph{arXiv:2601.12294}.

\bibitem[{Traj-RM}(2026)]{traj_reward2026}
Trajectory-level reward modeling.
2026.
\emph{arXiv:2604.08178}.

\end{thebibliography}

\end{document}
